% Section 3.2. Worked example: results. The claims table moved to Supplement S4.8 in the 4 Oct 2026 condense pass.

\subsection{Results}
\label{sec:results}

\begin{table}[!ht]
\centering
\caption{Ladder Verdicts per Model in the Worked Example}
\label{tab:verdicts}
\small
\newcolumntype{P}[1]{>{\raggedright\arraybackslash}p{#1}}
\begin{tabular}{@{}lP{2.1cm}P{2.3cm}P{1.9cm}P{1.9cm}P{2.5cm}@{}}
\toprule
Model & Gate: draws needed (min.\ / rec.) & Level 1: single draws & Level 2: income gap ratio & Level 3: overall residual & Level 4: structure \\
\midrule
DeepSeek-V3 & Pass (3 / 9) & Fail, compressed & Fail, 1.91 & Fail, +4.6 & Fail, no general factor \\ % R: gate_dstudy.csv; l2_headline.csv; 80d_headline.csv
Gemini-3-Flash & Pass (3 / 10) & Fail, compressed & Fail, 5.07 & Fail, +2.6 & Fail, loadings (R2) and invariance (R4) \\ % R: gate_dstudy.csv; l2_headline.csv; 80d_headline.csv
GPT-4o-mini & Pass (4 / 13) & Fail, compressed & Fail, 2.16 & Fail, +4.7 & Fail, no general factor \\ % R: gate_dstudy.csv; l2_headline.csv; 80d_headline.csv
GLM-4.7 & Pass (6 / 21) & Fail, too few atypical & Fail, 5.03 & Fail, +2.1 & Fail, loadings (R2) and invariance (R4) \\ % R: gate_dstudy.csv; l2_headline.csv; 80d_headline.csv
\bottomrule
\end{tabular}
\vspace{4pt}
\begin{tabnote}
\textit{Note.} Gate: draws needed for a persona mean's standard error to reach 0.25 (minimum) and 0.125 (recommended) NHANES standard deviations, the larger of the two framings. Level 2: ratio of the simulated to the standardized NHANES gap between low- and high-income personas, where 1 reproduces the gap. Level 3: post-stratified simulated minus NHANES mean PHQ-8, in points. Level 4: R2, loadings stronger than people's; R4, the factor differs across income where it does not in NHANES. Full per-level tables are in Supplement S4.
\end{tabnote}
\end{table}
\FloatBarrier

All four models pass the gate and fail every level above it (Table~\ref{tab:verdicts}), and Figure~\ref{fig:panels} shows what each failure looks like.

\begin{figure}[tp]
% Figure 2 title. Drawn by scripts/84e_fig2_ladder_panels.py from analysis/brm/gate_dstudy.csv,
% l1_summary.csv, 80d_headline.csv and l4_structure.csv. Every level-2 interval is in Supplement Figure S1.
\caption{Failures of the Four Models in the Worked Example}
\label{fig:panels}

{\centering\includegraphics[width=\linewidth]{fig2_ladder_panels.pdf}\par}
% Figure 2 note (below the figure): how to read each panel.
\figurenote{Gate: how far a persona's PHQ-8 total moves between draws, for one draw (the draw standard deviation, larger framing) and for the mean of 30 (its standard error), against the gate's limit of 0.98 points and recommended level of 0.49. % R: gate_dstudy.csv
Level 1: percent of draws more unusual (filled) or more regular (open) than 95\% of NHANES adults with the same total, both framings, where real people give 5\% of each.
Levels 2 and 3: post-stratified mean PHQ-8 total in each income band.
Level 4: loading of each PHQ-8 item on a single factor in the narrative framing, where a loading near 1 means the item rises and falls with the rest of the scale and a loading near 0 or below means it does not. DeepSeek-V3 and GPT-4o-mini have no general factor, and Gemini-3-Flash and GLM-4.7 follow the human pattern with loadings stronger than people's.
Every level-2 interval is in Supplement Figure S1.}

\end{figure}

\paragraph{Gate.}
No model can be read from a single draw, and two draws of one persona fall in different PHQ-8 severity bands 32.2\% to 35.2\% of the time across the two framings. % R: gate_bandflip_summary.csv
Averages of a few draws are stable enough to read, with the minimum rule needing 3 to 6 draws per persona. % R: gate_dstudy.csv
Setting the temperature to 0 narrows the spread of draws without separating personas any better (Supplement S1.7), and averaging remains the way to read a persona's score.

\paragraph{Level 1.}
No model passes level 1 in either framing, at any tolerance up to 3 or at temperature 0 (Supplement S4). % R: l1_personfit_verdicts_tau.csv; 96_decoding_l1.csv
Where 5\% of NHANES adults give an unusual answer pattern for their total, the models give 0.0\% to 1.0\%, and three of the four are compressed in both tails. % R: l1_summary.csv
A vignette bank or a screener test set drawn from these samples would miss almost all of the unusual presentations that one in twenty real respondents give.

\paragraph{Level 2.}
Every model fails level 2 on income, steepening the gap between low- and high-income personas to 1.91 to 5.07 times the population gap (Supplement Figure S1). % R: l2_headline.csv
We read the size of that ratio as partly a product of how the persona's income label is matched to the survey, although every model still fails on an income contrast when income is matched on the label's dollar amount (Supplement S3). % R: l2_verdicts.csv
On the sex gap the models disagree, two shrinking it, one doubling it and one unresolved, and pooled over the four models the ratio is 0.97, a near match built from errors in opposite directions. % R: l2_headline.csv
No racial contrast is kept, and NHANES measures the racial gaps among matched adults too loosely to certify any simulator. % R: l2_stop.csv; l2_headline.csv
None of these samples can estimate how depression differs by income, sex or race, and a study of the racial gaps needs a larger or more targeted reference before the ladder can read them.

\paragraph{Level 3.}
Every simulated population is more depressed than NHANES, by 2.1 to 4.7 points overall, and every group in every model sits above its real counterpart. % R: 80d_headline.csv; 80d_pass_counts.csv
The excess is largest in the low-income band, where 55.5\% of simulated respondents score 10 or more against 13.6\% of real adults, and it holds under all eight reference specifications (Supplement S3). % R: 80d_headline.csv; 89_l3_spec_range_summary.csv
Prevalence errs in a different pattern from the mean, and DeepSeek-V3's overall prevalence looks correct only because a low-income excess cancels shortfalls elsewhere. % R: 80d_model_verdicts.csv
Read as a stand-in for a survey, these samples would overstate depression for every group, most of all among low-income adults.

\paragraph{Level 4.}
DeepSeek-V3 and GPT-4o-mini have no general factor, and their PHQ-8 total cannot be scored as one measure. % R: l4_structure.csv
Gemini-3-Flash and GLM-4.7 have a factor with the human pattern of loadings, but loadings stronger than people's and a structure that changes across income where it holds in NHANES. % R: l4_structure.csv; l4_r2_size.csv; l4_invariance.csv
Their factor is built from differences between persona means, and draws of one persona share none, the same uniformity within personas that level 1 found (Supplement S4). % R: l4_between_within.csv

\paragraph{Summary.}
Read one at a time, every model's answers look like believable patients, and averages of a few draws are stable enough to read.
The introduction's persona shows the pattern on one case, with each of her draws plausible, the set less varied than people at her severity, and her mean well above matched NHANES adults in both framings. % R: 85_worked_case.csv
A researcher who checked this sample only by reading its answers or by the stability of its scores would have accepted it.
Checks of the kind used to report algorithmic fidelity would have looked reassuring as well.
Across the 48 matched persona cells the models' persona means correlate with the means of matching NHANES adults at .72 to .82, and every model gives at least five of the seven level-2 gaps the population's direction (Supplement S4.9). % R: 99_naive_baseline.csv
Neither a correlation nor a count of gap directions registers that every group is too depressed or that the income gradient is up to five times too steep.
The ladder fails it at every level above the gate, finding its answers too uniform, its income gradient too steep, every group too depressed and its total without the structure the PHQ-8 has in people.

\paragraph{Using the verdicts.}
We read the verdicts as a list of supported and ruled-out uses for a researcher holding this corpus.
Persona means can be read after a few draws, while single answers as cases, group comparisons, population estimates and scale totals are all ruled out for these samples.
The verdicts hold for these four models under this prompt family, and a new model, prompt or persona wording calls for the ladder to be run again.
A prompt control that asked each model to answer as the person described did not remove the excess at level 3, although it left unchanged the income labels, which paired a dollar amount with an insurance type (Supplement S1.6). % R: ../prompt_control_results.csv
